\section{A proved reduction at every even weight in the Gaussian ladder}
\label{gau:section}

The draft \emph{Gaussian multiple polylogarithm depth} studies
\begin{equation}\label{gau:sum-definition}
 S_p=\sum_{n\geq 0}\frac{(-1)^nH_n}{(2n+1)^p},
 \qquad H_n=\sum_{j=1}^n\frac1j,\quad H_0=0.
\end{equation}
Its formulas for $S_5$ and $S_7$ are correct, but the argument that a new
Dirichlet beta value ``absorbs'' the depth-two content is not a proof of a
reduction at every even weight. Here we supply a complete analytic proof and
an exact generating function.

\paragraph{Prior art and the precise contribution.}
The all-weight formula below already appears in a 2022 answer by
Ali Olaikhan, including exactly the $S_5$ and $S_7$ specializations displayed
in the repository draft \cite{gau-olaikhan-2022}.
It therefore requires attribution rather than a novelty claim. The purpose
of this section is to give a self-contained proof with convergence and
parameter domains specified, connect the formula to the draft's depth
question, and extract a convergent resummation. None of these arguments
asserts irrationality, algebraic independence, or minimal depth.

\subsection{Integral representation and a reflection identity}

For every integer $p\geq1$, the sum in \eqref{gau:sum-definition} converges.
For $p\geq2$ it converges absolutely. For $p=1$, the positive sequence
$H_n/(2n+1)$ decreases to zero: indeed,
\[
 \frac{H_{n+1}}{2n+3}\leq\frac{H_n}{2n+1}
 \quad\Longleftrightarrow\quad
 \frac{2n+1}{n+1}\leq2H_n,
\]
which holds for $n\geq1$. The alternating-series test applies.

\begin{lemma}\label{gau:moment}
For every integer $p\geq1$,
\begin{equation}\label{gau:moment-formula}
 S_p=-\frac1{(p-1)!}\int_0^1
       \frac{(-\log x)^{p-1}\log(1+x^2)}{1+x^2}\,dx.
\end{equation}
\end{lemma}
\begin{proof}
Summing the geometric series after the inner harmonic-number sum gives
\[
 \sum_{n\geq0}H_n y^n=-\frac{\log(1-y)}{1-y},\qquad |y|<1.
\]
First use $y=-r x^2$ with $0<r<1$, multiply by
$(-\log x)^{p-1}/(p-1)!$, and integrate term by term. The elementary
Mellin integral of $x^{2n}$ gives $(2n+1)^{-p}$.
Now let $r\uparrow1$. Abel's theorem applies to the convergent alternating
series. On the integral side,
$\log(1+rx^2)/(1+rx^2)$ is uniformly bounded on $0\leq x\leq1$,
and $(-\log x)^{p-1}$ is integrable, so dominated convergence applies.
This also covers the conditionally convergent case $p=1$.
\end{proof}

Introduce
\begin{align}
 J(s)&=\int_0^1\frac{x^{s-1}\log(1+x^2)}{1+x^2}\,dx,
 & C(s)&=\int_0^1\frac{x^{s-1}}{1+x^2}\,dx,\label{gau:mellin-defs}\\
 K(s)&=\int_0^\infty\frac{x^{s-1}\log(1+x^2)}{1+x^2}\,dx.
 \nonumber
\end{align}
The functions $J$ and $C$ are holomorphic for $\Re s>-2$ and $\Re s>0$,
respectively; $K$ is holomorphic for $-2<\Re s<2$.
In the common strip $0<\Re s<2$, the substitution $x=1/t$ in the
part of $K$ above $1$ gives
\begin{equation}\label{gau:reflection}
 K(s)=J(s)+J(2-s)-2C'(2-s).
\end{equation}
All signs in this identity matter: the logarithm transforms as
$\log(1+t^{-2})=\log(1+t^2)-2\log t$.

The beta integral gives another evaluation. For
$0<\Re s<2\Re a$,
\[
 \int_0^\infty\frac{x^{s-1}}{(1+x^2)^a}\,dx
 =\frac{\Gamma(s/2)\Gamma(a-s/2)}{2\Gamma(a)}.
\]
Taking minus the derivative in $a$ at $a=1$ yields
\begin{equation}\label{gau:complete-mellin}
 K(s)=\frac{\pi}{2\sin(\pi s/2)}
          \bigl[\psi(1)-\psi(1-s/2)\bigr],
 \qquad 0<\Re s<2.
\end{equation}
Differentiation under the integral is justified locally uniformly in that
strip by integrable power bounds with a logarithmic factor.

\subsection{The finite reduction}

Write
\[
 \beta(r)=\sum_{n\geq0}\frac{(-1)^n}{(2n+1)^r},
 \qquad
 \sec t=\sum_{j\geq0}\frac{|E_{2j}|}{(2j)!}t^{2j},
\]
so that $|E_0|,|E_2|,|E_4|,|E_6|,|E_8|=1,1,5,61,1385$.
Define $b_0=\log2$ and
$b_k=(1-2^{-2k-1})\zeta(2k+1)$ for $k\geq1$.

\begin{theorem}[All even-weight reductions]\label{gau:all-weights}
For every integer $m\geq1$,
\begin{equation}\label{gau:finite-reduction}
 \boxed{\displaystyle
 S_{2m-1}=(2m-1)\beta(2m)
 -\frac\pi2\sum_{k=0}^{m-1}
 \frac{|E_{2m-2-2k}|}{(2m-2-2k)!}
       \left(\frac\pi2\right)^{2m-2-2k}b_k.}
\end{equation}
Consequently $S_{2m-1}$ is a rational linear combination of
$\beta(2m)$, $\pi^{2m-1}\log2$, and
$\pi^{2m-2k-1}\zeta(2k+1)$, $1\leq k\leq m-1$.
\end{theorem}
\begin{proof}
For $|z|<1$, take the even part of \eqref{gau:reflection} about $s=1$.
If $K_{\mathrm{ev}}(z)=(K(1+z)+K(1-z))/2$, then
\begin{equation}\label{gau:even-reflection}
 J(1+z)+J(1-z)
 =K_{\mathrm{ev}}(z)+C'(1+z)+C'(1-z).
\end{equation}
The moment representation implies
\[
 J(1+z)+J(1-z)=-2\sum_{m\geq1}S_{2m-1}z^{2m-2}.
\]
Differentiating $C$ under its integral, then using the geometric series,
gives
\[
 C'(1+z)+C'(1-z)
 =-2\sum_{m\geq1}(2m-1)\beta(2m)z^{2m-2}.
\]
These expansions converge locally uniformly for $|z|<1$.

Next, \eqref{gau:complete-mellin} yields
\[
 K_{\mathrm{ev}}(z)=\frac\pi{2\cos(\pi z/2)}
 \left[\psi(1)-\frac{\psi((1-z)/2)+\psi((1+z)/2)}2\right].
\]
The digamma series, obtained by termwise differentiation of
$\psi(w)=-\gamma+\sum_{n\geq0}((n+1)^{-1}-(n+w)^{-1})$, gives
\[
 \psi(1)-\psi(\tfrac12)=2\log2,
 \qquad
 \psi^{(2k)}(\tfrac12)=-(2k)!(2^{2k+1}-1)\zeta(2k+1)
 \quad(k\geq1).
\]
It follows that
\[
 K_{\mathrm{ev}}(z)
 =\pi\sec(\pi z/2)\sum_{k\geq0}b_k z^{2k}.
\]
Substitute these three series into \eqref{gau:even-reflection} and compare
the coefficient of $z^{2m-2}$. The result is
\eqref{gau:finite-reduction}.
\end{proof}

For orientation, the first two instances are
\[
 S_1=\beta(2)-\frac\pi2\log2,
 \qquad
 S_3=3\beta(4)-\frac{\pi^3}{16}\log2
                     -\frac{7\pi}{16}\zeta(3).
\]
The cases $m=3,4$ give exactly the draft's equations
\texttt{eq:S5} and \texttt{eq:S7}. The next instance is
\begin{align}\label{gau:s9}
 S_9={}&9\beta(10)-\frac{277\pi^9}{4128768}\log2
       -\frac{427\pi^7}{737280}\zeta(3)
       -\frac{155\pi^5}{24576}\zeta(5)\nonumber\\
     &\hspace{14mm}-\frac{127\pi^3}{2048}\zeta(7)
       -\frac{511\pi}{1024}\zeta(9).
\end{align}
Every term has total weight $10$.

\paragraph{What parity does and does not prove.}
The mechanism is the reflection $s\mapsto2-s$ of the Mellin integral.
Its even Taylor part isolates the odd powers of the denominator. This
proves the upper bound on the complexity of $S_{2m-1}$ without assuming
that any even beta value is a new transcendental.
It gives no lower bound for the complexity of $S_{2m}$.
In particular, neither the collapse of $\beta(2m+1)$ to a rational multiple
of $\pi^{2m+1}$ nor a failed integer-relation search proves that
$S_{2m}$ requires a double polylogarithm.

\subsection{An exact generating function and a resummation}

\begin{proposition}\label{gau:generating}
The generating function
$F(z)=\sum_{m\geq1}S_{2m-1}z^{2m-2}$ has Taylor radius exactly $3$.
It is represented, for $|\Re z|<3$, by
\begin{equation}\label{gau:generating-integral}
 F(z)=-\int_0^1\frac{\cosh(z\log x)\log(1+x^2)}{1+x^2}\,dx.
\end{equation}
For $|z|<1$ it also has the explicit form
\begin{equation}\label{gau:generating-digamma}
\begin{split}
 F(z)={}&-\frac{C'(1+z)+C'(1-z)}2\\
 &-\frac{\pi}{4\cos(\pi z/2)}
 \left[\psi(1)-\frac{\psi((1-z)/2)+\psi((1+z)/2)}2\right],
\end{split}
\end{equation}
where
$C'(s)=\bigl[\psi'((s+2)/4)-\psi'(s/4)\bigr]/16$.
The apparent singularities of the whole right side at $z=\pm1$ are removable.
Moreover,
\begin{equation}\label{gau:resummation}
 \boxed{\displaystyle\sum_{m\geq1}S_{2m-1}=-\frac{\pi^2}{48}.}
\end{equation}
\end{proposition}
\begin{proof}
The alternating-series estimate gives $|S_p|\leq3^{-p}$ for $p\geq1$.
It follows that the Taylor series converges for $|z|<3$, and summing the
moments in \eqref{gau:moment-formula} gives
\eqref{gau:generating-integral}. The integrand is dominated on compact
subsets of $|\Re z|<3$ because $\log(1+x^2)/(1+x^2)=O(x^2)$ at zero.
Equation \eqref{gau:generating-digamma} follows from
\eqref{gau:even-reflection}. Its apparent singularities at $\pm1$ must
cancel by \eqref{gau:generating-integral}.

More explicitly, subtracting the constant coefficient gives the meromorphic
continuation
\begin{equation}\label{gau:partial-fractions}
 F(z)=S_1+\sum_{n\geq1}
 \frac{(-1)^nH_n z^2}{(2n+1)((2n+1)^2-z^2)}.
\end{equation}
The series on the right is absolutely and locally uniformly convergent
away from $z=\pm(2n+1)$, $n\geq1$. Its $n=1$ summand has a nonzero simple
pole at $z=3$, while all remaining summands are analytic there. Therefore
the Taylor radius is exactly $3$.
Finally, setting $z=1$ in \eqref{gau:generating-integral} simplifies the
kernel:
\[
 F(1)=-\frac12\int_0^1\frac{\log(1+x^2)}x\,dx
      =-\frac14\int_0^1\frac{\log(1+t)}t\,dt
      =-\frac{\pi^2}{48}.
\]
The last integral follows from the absolutely convergent integrated
alternating series for $\log(1+t)$.
\end{proof}

\paragraph{Reproducible checks.}
The program \texttt{code/verify\_gaussian.py} computes $p=1,3,5,7,9$ in
three ways: accelerated summation of \eqref{gau:sum-definition}, numerical
quadrature of \eqref{gau:moment-formula}, and the finite expression
\eqref{gau:finite-reduction}. At $90$ decimal digits the largest pairwise
absolute residual was less than $4\cdot10^{-91}$. It also checks five
values of the generating function, including a complex argument and
$z=3/2$, and the removable-singularity value $F(1)$.
The values and precision settings are recorded in
\texttt{data/gaussian\_checks.json}. These checks detect implementation
and transcription errors; the proofs above establish the identities.
